% SCP10, source lines 1818--1915, figs4/renorm-fixedpoint.pdf.
% Exact untwisted geometric blocking, including coarse periods one and two.

\begin{definition}[Four-site contraction with paired boundary legs]%
    \label{def:peps_two_by_two_tensor}
    \lean{TNLean.PEPS.twoByTwoBoundaryEquiv,
        TNLean.PEPS.twoByTwoSiteLegs,
        TNLean.PEPS.twoByTwoTensor,
        TNLean.PEPS.twoByTwoPeriodicBoundary}
    \leanok
    Let $X$ and $P$ be finite virtual and physical alphabets. Number the four
    corners clockwise from the upper right by $0,1,2,3$. Write
    $\alpha=(t,r,b,l)\in(X^2)^4$, with top and bottom pairs ordered left to
    right and right and left pairs ordered top to bottom. For
    $x=(x_0,x_1,x_2,x_3)\in X^4$, define the four fine leg configurations by
    \begin{align}
        \lambda_0(\alpha,x) & =(t_1,r_0,x_1,x_0),       &
        \lambda_1(\alpha,x) & =(x_1,r_1,b_1,x_2),\notag   \\
        \lambda_2(\alpha,x) & =(x_3,x_2,b_0,l_1),       &
        \lambda_3(\alpha,x) & =(t_0,x_0,x_3,l_0).
    \end{align}
    The eight separate boundary labels are indexed by
    $\{0,1,2,3\}\times\{0,1\}$. For tensors
    $a_i:X^4\times P\to\mathbb C$, their actual blocked tensor is
    \begin{align}
        B(\alpha,\sigma)=\sum_{x\in X^4}\prod_{i=0}^3
        a_i(\lambda_i(\alpha,x),\sigma_i).
    \end{align}
    On a periodic coarse lattice, with horizontal and vertical paired bond
    assignments $h,v$, the boundary at $z$ is
    $(v_z,h_z,v_{z-\hat y},h_{z-\hat x})$.
    This is the four-site grouping in
    \cite[Observation~6.6]{Schuch2010PEPS}.
\end{definition}

\begin{definition}[The actual periodic bond bijection]%
    \label{def:peps_two_by_two_bond_bijection}
    \lean{TNLean.PEPS.twoByTwoTiledBondEquiv,
        TNLean.PEPS.twoByTwoSiteLegs_periodic}
    \leanok
    \uses{def:peps_two_by_two_tensor}
    Let $C=\mathbb Z_w\times\mathbb Z_h$ with $w,h>0$, and identify the fine
    torus $\mathbb Z_{2w}\times\mathbb Z_{2h}$ with $C\times\{0,1,2,3\}$
    by its disjoint four-site tiling. The paired crossing bonds $h_z,v_z$
    and the four internal bonds $x_{z,i}$ give the fine outgoing bonds by
    \begin{align}
        H_z & =((h_z)_0,(h_z)_1,x_{z,2},x_{z,0}), \\
        V_z & =((v_z)_1,x_{z,1},x_{z,3},(v_z)_0).
    \end{align}
    This is a bijection. Reading the top, right, bottom, and left fine
    bonds gives exactly $\lambda_i(\alpha_z,x_z)$, including periodic seams.
    There is one horizontal and one vertical bond per site even when a
    coarse period is one or two.
\end{definition}

\begin{theorem}[Exact periodic four-site reblocking]%
    \label{thm:peps_two_by_two_geometric_blocking}
    \lean{TNLean.PEPS.torusBondNetwork_eq_twoByTwoTiledSum,
        TNLean.PEPS.torusBondNetwork_eq_twoByTwoBlocked}
    \leanok
    \uses{def:peps_two_by_two_bond_bijection}
    Place arbitrary, possibly different tensors $a_z$ on the fine torus,
    and let $B_v$ be their actual four-site contractions. With identity
    operators on every bond, the original fine state satisfies
    $\Psi_a(\sigma)=\Psi_B(\tau)$, where
    $\tau_{v,i}=\sigma_{(v,i)}$. Both sides sum over all their bonds.
    The equality has scalar factor one and holds for every $w,h>0$.
\end{theorem}

\begin{proof}\leanok
    \uses{def:peps_two_by_two_bond_bijection}
    Reindex the fine horizontal and vertical sums by the displayed
    bijection. The remaining internal sum separates as
    $\sum_{x:C\to X^4}\prod_v f_v(x_v)=\prod_v\sum_{x_v\in X^4}f_v(x_v)$.
    Its factor at $v$ is precisely $B_v(\alpha_v,\tau_v)$.
\end{proof}

\begin{definition}[Physical four-site regrouping]%
    \label{def:peps_two_by_two_physical_grouping}
    \lean{TNLean.PEPS.twoByTwoPhysicalEquiv,
        TNLean.PEPS.twoByTwoPhysicalMatrix}
    \leanok
    The disjoint tiling identifies fine physical configurations $\sigma$
    with block configurations $\tau_{v,i}=\sigma_{(v,i)}$.
    Let $W$ be its permutation matrix on the complete physical spaces.
\end{definition}

\begin{theorem}[Unitary physical regrouping of the fine state]%
    \label{thm:peps_two_by_two_physical_grouping}
    \lean{TNLean.PEPS.twoByTwoPhysicalMatrix_isIsometry,
        TNLean.PEPS.twoByTwoPhysicalMatrix_conjTranspose_isIsometry,
        TNLean.PEPS.twoByTwoPhysicalMatrix_mulVec_network}
    \leanok
    \uses{def:peps_two_by_two_physical_grouping,
        thm:peps_two_by_two_geometric_blocking}
    Both $W$ and $W^*$ are isometries, and $W\Psi_a=\Psi_B$.
    All four physical registers per block, including unused directions,
    belong to this unitary identification.
\end{theorem}

\begin{proof}\leanok
    The physical label map is bijective, so its matrix and the inverse
    permutation matrix preserve inner products. Evaluating its action on
    each coefficient gives the exact reblocking equality.
\end{proof}

\begin{definition}[Incident coordinates of the coarse block]%
    \label{def:peps_two_by_two_graph_site}
    \lean{TNLean.PEPS.torusIncidentFamily,
        TNLean.PEPS.twoByTwoGraphSite}
    \leanok
    \uses{def:peps_two_by_two_tensor}
    Suppose $w,h\geq3$, so the four incident edges at every coarse vertex
    are distinct. A four-leg tensor is expressed on incident labels by
    reading its top, right, bottom, and left edges. Applying this change
    to $B_v$ defines its coarse graph tensor, with physical alphabet $P^4$.
\end{definition}

\begin{theorem}[Equality with the actual coarse graph contraction]%
    \label{thm:peps_two_by_two_graph_contraction}
    \lean{TNLean.PEPS.graphBondNetwork_torusIncidentFamily,
        TNLean.PEPS.torusBondNetwork_eq_twoByTwoGraphSite}
    \leanok
    \uses{def:peps_two_by_two_graph_site,thm:peps_two_by_two_geometric_blocking}
    For $w,h\geq3$, the bond-indexed contraction equals the incident-edge
    graph contraction, for arbitrary site-dependent tensors and physical
    alphabets. In particular $W\Psi_a$ is the actual graph state of the
    paired four-site tensors.
\end{theorem}

\begin{proof}\leanok
    Every coarse graph edge is uniquely a horizontal or vertical outgoing
    edge. Reindexing its label sum yields the bond-indexed contraction.
    Apply the exact fine-to-coarse reblocking equality.
\end{proof}
